% Demonstration layout: keep each evidence graphic with its interpretation.
% A native minipage caption avoids float-only pages without shrinking the figure.
\makeatletter
\newenvironment{demofixedfigure}{%
  \par\medskip\noindent\begin{minipage}{\linewidth}\def\@captype{figure}\centering
}{\end{minipage}\par\medskip}
\makeatother
\section{Problem Restatement and Analysis}
The demonstration asks three linked questions: how closely two numerical
methods reproduce a known solution, how the solution changes with the decay
rate, and how an assumed range of rates propagates into an output envelope.
Every value is synthetic. Figure~\ref{fig:demo-method} connects the model,
analytical reference and checks; it is a method diagram rather than evidence
about a physical system.
\begin{demofixedfigure}
\centering\includegraphics[width=6.5in]{figure-demo/method.pdf}
\caption{Synthetic method overview: an assumed decay equation, two numerical
solvers, an analytical reference and numerical comparison.}
\label{fig:demo-method}
\end{demofixedfigure}

\section{Assumptions and Symbols}
The quantity is dimensionless, its initial value is one, and the rate is
constant within each scenario. These are deliberate simplifying assumptions,
not measurements. Time is denoted by $t$, the remaining fraction by $y(t)$, the
positive rate by $k$, and the fixed numerical step by $h$.

\clearpage
\section{Model, Solution and Validation}
The shared model and its exact comparison reference are
\begin{equation}
\begin{aligned}
\frac{\mathrm{d}y}{\mathrm{d}t}&=-k y, & y(0)&=1,\\
y(t)&=\exp(-k t), & k&>0.
\end{aligned}\label{eq:demo-decay}
\end{equation}
Euler uses $y_{n+1}=y_n-hky_n$; Runge--Kutta combines four intermediate slopes.
Both use $h=\FigureDemoStep$, $k=1$, and 41 shared time points in $[0,4]$.
We compare both numerical solutions with Equation~\ref{eq:demo-decay}.

\subsection{Results and interpretation}
The largest absolute grid errors are \FigureDemoEulerError\ for Euler and
\FigureDemoRKError\ for Runge--Kutta. Figure~\ref{fig:demo-comparison}
puts both methods on the same comparison scale.
\begin{demofixedfigure}
\centering\includegraphics[width=6.5in]{figure-demo/comparison.pdf}
\caption{Synthetic numerical comparison. All lines use the same 41 time
points; Runge--Kutta markers are shown every fourth point for clarity.
Blue: analytical reference; dashed orange: Euler; dotted green: Runge--Kutta.
Values and units are dimensionless.}
\label{fig:demo-comparison}
\end{demofixedfigure}
Euler decays too quickly on this grid, while Runge--Kutta follows the exact
reference closely. This supports accuracy for this smooth test equation and
step; it does not establish accuracy for another model or time scale.

\clearpage
\section{Sensitivity and Uncertainty Analysis}
\subsection{Parameter sensitivity}
Figure~\ref{fig:demo-sensitivity} compares rates of $0.7$, $1.0$ and $1.3$
with identical initial conditions. Increasing the rate lowers the remaining
fraction at every positive time. These are scenario responses rather than
fits to observations.
\begin{demofixedfigure}
\centering\includegraphics[width=6.5in]{figure-demo/sensitivity.pdf}
\caption{Synthetic sensitivity to the assumed constant decay rate.
Line style supplements color so the three scenarios remain distinguishable.}
\label{fig:demo-sensitivity}
\end{demofixedfigure}
The shared initial value isolates the effect of the rate. The curves separate
over time, showing why a model-based decision should report its assumed rate
rather than present a single decay curve without qualification.

\clearpage
\subsection{A parameter envelope}
We evaluate \FigureDemoRateCount\ equally spaced rates in $[0.7,1.3]$ and
retain all scenarios. Figure~\ref{fig:demo-uncertainty} shows the pointwise
10th--90th percentile envelope of those specified rates. It describes this
chosen parameter grid, not repeated observations, a fitted posterior or a
sampling confidence interval.
\begin{demofixedfigure}
\centering\includegraphics[width=6.5in]{figure-demo/uncertainty.pdf}
\caption{Synthetic propagation of a specified parameter range. Shading is
the pointwise 10th--90th percentile across the 301 deterministic rate
scenarios; the solid line uses $k=1$. The band is not a confidence interval.}
\label{fig:demo-uncertainty}
\end{demofixedfigure}

\section{Model Evaluation and Conclusions}
The analytical reference, shared grid and exported data make the numerical
comparison easy to check. The four graphics have distinct roles: explain
the method, compare solvers, assess sensitivity and display uncertainty.
The limitations are equally important: there are no observations, no
parameter estimation and no independent agent's scientific attestation.
Nonconstant rates and noisy measurements would require a new model and
appropriate validation. This example demonstrates the reporting pipeline,
not the validity of a real-world prediction. Contest presentation and AI
disclosure still follow the applicable rules~\cite{comap2027rules,comap2025ai}.
